\begin{abstract}
Novel view synthesis from dynamic streaming video faces a fundamental trade-off: maintaining a persistent, long-horizon memory to reconstruct temporarily occluded regions while operating under strict real-time constraints. While Test-Time Training (TTT) offers a powerful memory mechanism, standard models mandate gradient-based memory updates at every frame to adapt to the changing motion in dynamic scenes. This computational overhead from heavy memory updates precludes real-time application and can lead to instability over extended contexts. Recognizing that memory updates are considerably more demanding than memory application and that video content is largely redundant, we propose to decouple the frequencies of these two process. Our approach performs periodic memory updates while applying the memory on a per-frame basis, utilizing cross-view attention to manage deformations between the prior memory state and the current frame. To lock in the historical context, we introduce two critical mechanisms: an auxiliary \emph{Memory Loss} that forces persistent internalization of the scene, and a \emph{Memory Caching} strategy that regularizes active weights against catastrophic drift. Extensive evaluations on unconstrained, dynamic human motion demonstrate that our method exhibits state-of-the-art performance, and successfully demonstrates minute-scale online memorization while operating in real-time.

% Novel view synthesis from streaming video inherently faces a fundamental conflict: memorizing over a long-context to groundedly inpaint invisible regions while maintaining real-time performance. While test-time training offers a powerful mechanism for memory modules, its significant computational overhead from heavy memory updates often precludes real-time application and can lead to instability over extended contexts. Recognizing that memory updates are considerably more demanding than memory application and that video content is largely redundant, we propose a novel framework that disentangles these two processes. Our approach performs periodic memory updates while applying the memory on a per-frame basis, utilizing a transformer to manage deformations between the prior memory state and the current frame. To ensure robust memorization under sparse disocclusion supervision, we introduce a memory auxiliary loss. Our method achieves long minute-scale online memorization capabilities while operating in real-time. We evaluate our framework across both human and large-scale in-the-wild scenes and demonstrate competitive or state-of-the-art performance.

\end{abstract}